\section{Introduction}
\label{sec:introduction}

Mixture-of-experts (MoE) models route each token to a small subset of a
much larger set of expert sub-networks~\cite{shazeer2017sparsely},
decoupling total parameter count from per-token compute. When the full
expert set does not fit resident in fast memory, inactive experts must
be retained, fetched, and evicted under a dynamically changing,
token-by-token routing workload -- a problem a growing line of
memory-constrained MoE inference systems addresses by offloading,
caching, or predicting expert demand~\cite{eliseev2023offloading,xue2025moeinfinity,kamahori2025fiddler}. The quality of the retention
policy governing this process directly determines how much expert data
must be re-fetched during inference.

This report presents a cumulative tranche of a three-stage research
programme (``Research Part~2''), organized around one central question:
\emph{how much of the cost of imperfect expert retention is an
irreducible property of the workload, and how much is recoverable by a
better, practical retention policy?} The programme is structured as
\begin{enumerate}
 \item \textbf{Research II-A -- Oracle headroom.} Establish, under a
   validated offline replay model, the \emph{upper bound} on recoverable
   retention performance: the gap between realistic causal policies and a
   non-causal, information-unbounded oracle.
 \item \textbf{Research II-B -- Routing structure.} Characterize whether
   the expert-routing process exhibits exploitable temporal or
   statistical structure (frequency concentration, transition structure,
   reuse-distance/locality structure, and workload dependence/drift) that
   could, in principle, explain and motivate a causal predictor.
 \item \textbf{Research II-C -- Causal recovery.} Build an actual causal
   predictor and measure what fraction of the II-A upper bound it
   recovers, including validation against the real runtime.
\end{enumerate}

This document reports the completed II-A and II-B stages.
Research~II-C has not started (Section~\ref{sec:iic-placeholder} states
its intended scope as a placeholder only); no result, finding, or claim
from II-C is made here. The II-A deliverable is the
\emph{policy-information headroom} between causal retention policies and a
non-causal oracle, decomposed from the \emph{structural} cost of the
runtime's sharded retention architecture, measured across a fixed sweep of
expert-memory budgets on a real, validated routing trace. The II-B
deliverable is a descriptive characterization of routing concentration,
same-layer transition dependence, short-horizon reuse locality, and
workload dependence under the tested protocol. II-A headroom is a
theoretical upper bound on recoverable performance, not demonstrated
causal recoverability; II-B structure is evidence that causal prediction
is scientifically worth testing, not evidence that prediction works or
improves runtime. That distinction is the central methodological concern
of this cumulative tranche.
